% ===================================================================
%  GAMBAR No. 8 — Panen. --- Mbebedhag, Panen anggur, Mancing.
%
%  Traduction, faite en 2026, en javanais d'aujourd'hui, sur l'ido de
%  texto/io. Regles de la colonne : voir 10-gambar-01.tex. Deux
%  scenes, deux numerotations.
%
%  UN ECART DECLARE DANS CE TABLEAU, ET IL REND SON BLOC AVEUGLE. Le
%  fac-simile ido compose « gobleto \textsuperscript{13)} » sans la
%  parenthese ouvrante, au bloc t08-c2-03-1 ; renvoji.py ne peut donc
%  pas lire ce renvoi-la, le bloc figure parmi les trois ecarts
%  declares, et le controle n'y regarde PLUS RIEN — pas seulement la
%  parenthese en cause. L'ordre 3, 12, 14, 13, 15, 16 a donc ete pose
%  et relu a la main, sur le francais qui l'imprime bien. C'est
%  exactement la faute que la colonne tamoule avait faite ici meme :
%  elle avait recopie « 13) » du fac-simile en croyant le controle
%  derriere elle. Une traduction n'est pas une transcription, et un
%  ecart declare est l'endroit ou il faut le plus se relire.
%
%  LES DEUX MOISSONS DE CE TABLEAU N'EXISTENT PAS A JAVA, ET LA
%  COLONNE LES ECRIT AVEC LES MOTS DE JAVA QUAND ILS EXISTENT. Le
%  froment ne pousse pas ici et se dit « gandum », mot arabe venu par
%  le malais ; mais la moisson elle-meme se dit « panen », et
%  « ngunduh » et « ngarit » sont les gestes javanais qu'on fait sur le
%  riz depuis toujours. La vigne, elle, n'a rien : anggur est persan,
%  arrive par le malais, et tout ce qui l'entoure est emprunte. Une
%  langue prete ses VERBES a une plante qu'elle ne connait pas ; elle
%  ne lui prete pas ses NOMS.
%
%  ET LE MOULIN A VENT (50) EST LE CONTRAIRE DU MOULIN A EAU DU
%  TABLEAU 7. Java n'a jamais eu de moulin a vent — pas de vent
%  regulier, et l'eau partout —, et le javanais dit donc « kincir
%  angin », roue-a-vent, un compose transparent fait de deux mots
%  javanais pour une chose qui n'est pas venue. C'est le meme procede
%  que « sepatu es » au tableau 7 : quand la chose manque, la langue
%  decrit au lieu d'emprunter, et c'est le signe qu'elle n'a jamais eu
%  besoin d'en parler.
%
%  L'ARC-EN-CIEL (52) EST LE PLUS BEAU MOT DU TABLEAU, ET IL EST
%  JAVANAIS SANS PARTAGE : « kluwung ». L'indonesien dit « pelangi »,
%  qui est un autre mot malais ; le francais dit un arc dans le ciel,
%  l'ido copie le francais, l'anglais dit un arc de pluie. Aucun des
%  quatre ne dit ce que dit le javanais, qui ne le decompose pas du
%  tout : kluwung est la chose, sans image et sans arc.
%
%  ET LE TONNERRE ET LA FOUDRE SE SEPARENT ICI COMME ILS SE SEPARENT
%  RAREMENT : « gludhug » le grondement, « bledheg » l'eclair qui
%  tombe, « kilat » la lueur. Trois mots pour ce que le francais
%  partage en deux et l'ido en deux aussi. Un ciel qu'on regarde tous
%  les jours se dit plus finement qu'un ciel qu'on traduit.
% ===================================================================

%%K t08-noto-1 noto
\textit{(*) Amarga tembung iku ora cetha : apa iku panen lena, panen
anggur, apa panen apel? Bisa uga becik nganggo «} \VUgras{mesono}
\textit{» kang diusulake ing Mondo XIV, kaca 242. Nanging nganti saiki
oyod anyar iku durung ditampa dening Akademi lan isih ngenteni
rembugan miturut pranatan Idis kang lumaku.}

%%K t08-apar-1 apar
\VUsaut{5.0mm}
\VUtitre{15.0pt}{60}{GAMBAR No. 8}
\VUsaut{3.0mm}
\VUcentre{13.2pt}{90}{\textit{Adegan Kapisan.}}
\VUsaut{3.0mm}
\VUpk{10.2pt}{40}{Panen (*).}
\VUsaut{4.0mm}
\VUpk{10.2pt}{40}{Kaanane Pedesan.}
\VUsaut{3.0mm}

%%K t08-c1-01-1 p
1. --- Bareng karo \VUgras{mangsa panas,} kaanane pedesan owah maneh.
Wiji kang dipasrahake marang jugangan wis thukul ; tanduran cilik ijo
metu saka lemah lan wis awoh wuli. Wiji wis kadadeyan, banjur mateng,
lan saiki kari ngunduh wae.

%%K t08-c1-02-1 p
2. --- \VUgras{Kembang pedesan} wis padha megar. \VUgras{Kembang
bluwet} \textsuperscript{(1)}, \VUgras{kembang menur abang}
\textsuperscript{(2)} lan \VUgras{kembang jari} \textsuperscript{(3)}
nyampur ing warna kang sarwa padha ing sawah gandum, dadi bedane kang
bungah saka \VUgras{kelopak-kelopake} kang seger.

%%K t08-c1-03-1 p
3. --- Wit woh-wohan wis padha ngrembuyung godhonge ; wohe wis mateng.
\VUgras{Wit ceri} \textsuperscript{(4)} cilik kebak \VUgras{woh ceri}
\textsuperscript{(5)} endah kang diunduh lan dipangan bocah-bocah kanthi
seneng banget.

%%K t08-c1-04-1 p
4. --- Saiki kewan-kewan bisa nglampahi sedina muput ing jaba.

%%K t08-c1-04-2 p
\VUgras{Cempe} \textsuperscript{(6)} wis padha gedhe lan dadi
\VUgras{wedhus wadon} \textsuperscript{(7)} endah lan \VUgras{wedhus
lanang} \textsuperscript{(8)} endah kang wulune kandel. Kewan-kewan iku
kanthi tentrem nyenggut \VUgras{suket} ing \VUgras{pangonan}
\textsuperscript{(9)} kang kembang-kembang \VUgras{margrit}e.

%%K t08-c1-04-3 p
Sedhela maneh wong bakal wiwit nyukur \VUgras{kewan} iku supaya oleh
\VUgras{wul}e, kang digawe lawon sandhangan kita.

%%K t08-tit-2 sub
\VUsaut{5.0mm}
\VUpk{10.2pt}{40}{Pangon.}
\VUsaut{3.4mm}

%%K t08-c1-05-1 p
5. --- \VUgras{Pangon} \textsuperscript{(10)} katon ora patiya nggagas
kewane. Kang dipikir mung prahara kang liwat ing wates langit, dene
\VUgras{tukang jaga wedhus} \textsuperscript{(13)}, kang nyedhakake
\VUgras{kendhi banyune} \textsuperscript{(14)} ing lambene, katon luwih
ora ambil pusing maneh.

%%K t08-c1-06-1 p
6. --- Panase nyenyet ; amarga \VUgras{asu} \textsuperscript{(15)} uga
teka ngilangi ngelake ; asu iku ndilat banyu seger ing \VUgras{kali
cilik} \textsuperscript{(16)} lan gawe kagete \VUgras{kodhok cilik}
\textsuperscript{(17)} kang mesakake, kang ndhelik ing suket pinggir
kali. Kodhok iku mancolot menyang banyu bening panggonan
\VUgras{tunjung} \textsuperscript{(18)} kembang.

%%K t08-c1-07-1 p
7. --- \VUgras{Manuk kicuit} \textsuperscript{(19)} adus ing cedhake
\VUgras{sumber cilik} \textsuperscript{(20)} kang muncrat ing antarane
watu loro, banjur ndhelik ing sangisore \VUgras{grumbulan eri}
\textsuperscript{(21)} lan \VUgras{grumbulan krateg}
\textsuperscript{(22)}.

%%K t08-tit-3 sub
\VUsaut{5.0mm}
\VUpk{10.2pt}{40}{Panen.}
\VUsaut{3.4mm}

%%K t08-c1-08-1 p
8. --- Kanggo bocah-bocah desa, panen gandum iku wektu kang nyenengake
banget. Bocah-bocah kanthi bungah \VUgras{jempalikan}
\textsuperscript{(24)} ing \VUgras{tumpukan damen} \textsuperscript{(23)},
padha ngewangi \VUgras{wong panen} \textsuperscript{(26)}, lan nalika
wong-wong iku ngarit \VUgras{sawah gandum} \textsuperscript{(27)}
nganggo \VUgras{arite} \textsuperscript{(25)} kang dilandhepake becik
nganggo \VUgras{watu asahan} \textsuperscript{(30)}, bocah-bocah
nglumpukake gandum lan mbentheli dadi \VUgras{pocongan}
\textsuperscript{(31)} utawa \VUgras{bongkokan} \textsuperscript{(32)}.

%%K t08-c1-09-1 p
9. --- Kadhang bocah kang paling gedhe diparengake ngethok nganggo
\VUgras{arit bengkong} \textsuperscript{(12)} \VUgras{batang emas}
\textsuperscript{(28)} kang nyangga \VUgras{wuli} \textsuperscript{(29)}
kang abot kebak wiji ; dene bocah kang cilik-cilik, kang ngawasi
\VUgras{kewan} \textsuperscript{(33)} kang nyenggut ing \VUgras{pangonan}
\textsuperscript{(34)} ing sangisore ayange \VUgras{wit linden}
\textsuperscript{(35)} kang gedhe, nyekel \VUgras{kupu} kang endah
nganggo \VUgras{jaring}e \textsuperscript{(36)}.

%%K t08-c1-09-2 p
Malah sawetara menek ing \VUgras{gerobag} \textsuperscript{(37)} kanggo
nata pocongan kang diulungake marang bocah-bocah nganggo
\VUgras{garpu} \textsuperscript{(38)}.

%%K t08-c1-10-1 p
10. --- Ing saindenging \VUgras{dhataran} \textsuperscript{(39)},
panggonan \VUgras{manuk branjangan} \textsuperscript{(41)} wiwit mabur,
kabeh katon rame. Kabeh wong repot ngurusi panen. Nanging, nalika wong
mlarat isih nganggo piranti lawas — \VUgras{arit, arit bengkong} lan
\VUgras{gebugan} \textsuperscript{(40)} —, juragan sugih ngaritake gandum
utawa \VUgras{gandum ireng}e \textsuperscript{(43)} nganggo \VUgras{mesin
arit} \textsuperscript{(42)}. Banjur, sarehne ora perlu nggawa pocongan
menyang lumbung, wong nggawe \VUgras{tumpukan pocongan}
\textsuperscript{(44)} kang gedhe.

%%K t08-c1-11-1 p
11. --- Ing kono wong nekakake \VUgras{mesin perontok}
\textsuperscript{(47)} kang digerakake \VUgras{lokomobil}
\textsuperscript{(45)} nganggo \VUgras{sabuk transmisi}
\textsuperscript{(46)}, lan mangkono wiji dipisah saka \VUgras{damen}
tanpa ninggal sawah panggonane disebar.

%%K t08-tit-4 sub
\VUsaut{5.0mm}
\VUpk{10.2pt}{40}{Prahara.}
\VUsaut{3.4mm}

%%K t08-c1-12-1 p
12. --- Kerep \VUgras{prahara gedhe} \textsuperscript{(53)} kedadeyan ing
mangsa panen. Kang kapisan gludhug krungu saka kadohan ;
\VUgras{angin gedhe saka kulon} \textsuperscript{(54)} wiwit sumribit lan
mbengkongake wit kang paling gedhe ing \VUgras{alas}
\textsuperscript{(49)}. \VUgras{Mendhung} \textsuperscript{(56)} ireng
nyerbu langit ; mendhung iku ditembus kilatan \VUgras{bledheg}
\textsuperscript{(55)} kang nyilaukake. \VUgras{Udan} lan kadhang
\VUgras{udan es} wiwit tiba, ngremuk panenan kang wis siyap diarit.

%%K t08-c1-13-1 p
13. --- Kerep bledheg ngobong bangunan. Kaya taun kepungkur payone
\VUgras{kincir angin} \textsuperscript{(50)} dadi awu lan \VUgras{elare}
\textsuperscript{(51)} loro remuk.

%%K t08-c1-14-1 p
14. --- Sawise sawetara jam, kaanan tentrem maneh. \VUgras{Kluwung}
\textsuperscript{(52)} kang mencorong katon ing wates langit lan alam
mbaleni uripe kang kapedhot sawetara wektu. Wong desa, kang ngungsi
ing \VUgras{gubug} \textsuperscript{(48)}, wiwit nyambut gawe maneh lan
kanthi kendel ngudi ndandani karusakan kang lagi wae dialami.

%%K t08-tit-5 sub
\VUsaut{6.0mm}
\VUcentre{13.2pt}{90}{\textit{Adegan Kapindho.}}
\VUsaut{3.6mm}

%%K t08-tit-6 sub
\VUpk{10.2pt}{40}{Mbebedhag. --- Panen anggur.}
\VUsaut{1.4mm}
\VUpk{10.2pt}{40}{Mancing.}
\VUsaut{4.0mm}

%%K t08-c2-01-1 p
1. --- Ing sakubenge pungkasan sasi \VUgras{Agustus} utawa tengahe
\VUgras{September,} gumantung tlatahe, mangsa mbebedhag diwiwiti. Lagi
wae sunar srengenge kapisan gawe \VUgras{kadhal} \textsuperscript{(2)}
metu saka rongé, \VUgras{wong mbebedhag} \textsuperscript{(5)} wis
mangkat karo \VUgras{asu-asune} \textsuperscript{(4)}.

%%K t08-c2-02-1 p
2. --- Dheweke nganggo \VUgras{sandhangan mbebedhag}
\textsuperscript{(9)} kang kuwat saka lawon kandel, lan nganggo
\VUgras{sarung sikil} kulit, kang gawe dheweke bisa mlaku ing suket
teles panggonan \VUgras{kodhok} \textsuperscript{(1)} ndhelik ; dheweke
njupuk \VUgras{bedhil}e \textsuperscript{(6)} lan \VUgras{tas buron}e
\textsuperscript{(10)}, lan iki dheweke mangkat.

%%K t08-c2-03-1 p
3. --- Semangate anggone mburu nggawa dheweke menyang tengahe kebon
anggur panggonan panen anggur lagi rame banget. Anak-anake tukang
kebon anggur, kang lumrahe njaga wedhus jawa ing dalan, dina iki
nglilakake wedhus iku nyenggut sasenenge, lan sawise nalekake ing
patok \VUgras{wedhus jawa lanang} \textsuperscript{(3)} tuwa kang mesthi
bakal nganggo kebebasane kanggo mlayu, bocah-bocah iku uga melu
seneng-seneng. Kang siji njupuk dhompolan anggur saka \VUgras{kranjang
panen} \textsuperscript{(12)} lan seneng-seneng ngilekake \VUgras{sarine}
\textsuperscript{(14)} menyang \VUgras{cangkir} \textsuperscript{(13)}
kanggo nggawe anggur legi kang durung peragian. Sijine masang
\VUgras{makutha godhong} \textsuperscript{(15)} kang lagi wae dianyam.

Ing cedhake bocah-bocah iku, \VUgras{manuk emprit}
\textsuperscript{(16)} kang ora duwe isin teka nganti ing ngarepe alat
peres kanggo nyolong woh anggur.

%%K t08-c2-04-1 p
4. --- Nalika bocah-bocah seneng-seneng, wong lanang padha nyambut
gawe. \VUgras{Wong metik anggur} \textsuperscript{(18)} nggawa woh anggur
menyang kamar ngisor lemah nganggo \VUgras{kranjang gendhong}
\textsuperscript{(17)} ; \VUgras{tukang tong} \textsuperscript{(19)}
ndandani \VUgras{tong} \textsuperscript{(20)}, mriksa apa
\VUgras{papan sisih}e \textsuperscript{(22)} lan \VUgras{bunthut}e
\textsuperscript{(21)} isih becik, lan ngganti \VUgras{simpai}
\textsuperscript{(23)} kang pedhot. Dheweke uga wis nggarap \VUgras{bak
gedhe} \textsuperscript{(25)} panggonan wong ngesokake \VUgras{woh
anggur} \textsuperscript{(26)} supaya peragian.

%%K t08-c2-05-1 p
5. --- Yen peragiane wis rampung, wong nyaring anggure lan nyelehake
ampase ing \VUgras{alat peres} \textsuperscript{(28)}, banjur kanthi
alon-alon ngencengake \VUgras{sekrup gedhe} \textsuperscript{(29)} wong
meres sarine, kang mili menyang \VUgras{bak cilik}
\textsuperscript{(32)}, lan isih oleh \VUgras{anggur}
\textsuperscript{(31)} maneh.

%%K t08-tit-8 sub
\VUsaut{6.0mm}
\VUpk{10.2pt}{40}{Bir lan Sider.}
\VUsaut{3.4mm}

%%K t08-c2-06-1 p
6. --- Ing temboke \VUgras{kamar ngisor lemah} \textsuperscript{(27)}
mrambat pange \VUgras{hop} \textsuperscript{(30)}. Ing kono iku mung
tanduran rerenggan, nanging ing sawetara negara, panggonan anggur arang
banget, hop iku ditandur kanggo nggawe \VUgras{bir.}

%%K t08-c2-07-1 p
7. --- \VUgras{Wit apel} \textsuperscript{(33)} cilik kebak apel kang
yen wis mateng bakal ngasilake \VUgras{sider} kang enak banget. Ing
\VUgras{kurungane} \textsuperscript{(34)}, \VUgras{manuk kenari}
\textsuperscript{(35)} lan \VUgras{manuk kardhuel} \textsuperscript{(36)}
nyawang kanthi mripat mesgul marang manuk emprit kang bebas pethakilan.

%%K t08-c2-08-1 p
8. --- Nganti tekan wates pandeleng, ing ereng-erenge punthuk, jejer
\VUgras{cagak anggur} \textsuperscript{(40)} kang nyangga \VUgras{watang
anggur} \textsuperscript{(39)} kang endah banget, kebak \VUgras{dhompolan}
kang abot.

%%K t08-c2-08-2 p
Sataun muput, \VUgras{tukang kebon anggur} \textsuperscript{(42)} wis
nyambut gawe ngopeni \VUgras{kebon anggure} \textsuperscript{(38)}, lan
saiki dheweke diganjar panen kang endah kanggo rekasane. \VUgras{Wong
wadon kang metik anggur} \textsuperscript{(41)} ngiseni bak kang
diselehake ing gerobag kang ditarik \VUgras{kuldi}
\textsuperscript{(43)}, lan kabeh padha bungah.

%%K t08-tit-9 sub
\VUsaut{5.0mm}
\VUpk{10.2pt}{40}{Mancing.}
\VUsaut{3.4mm}

%%K t08-c2-09-1 p
9. --- Semono uga \VUgras{wong mancing} \textsuperscript{(44)} kang wiwit
esuk wis teka manggon ing pinggire banyu karo \VUgras{walesan}e
\textsuperscript{(45)}.

%%K t08-c2-09-2 p
\VUgras{Iwak} \textsuperscript{(47)} nyamber \VUgras{pancing} sanalika
wong mancing nyemplungake \VUgras{benang}e \textsuperscript{(46)} ing
banyu, lan lagi wae wong mancing sempat ngganti \VUgras{umpan}e,
korban anyar wis kesangkut ing pucuking benang.

%%K t08-c2-09-3 p
Nanging \VUgras{wong njala} \textsuperscript{(48)} kang, saka
\VUgras{prau}ne \textsuperscript{(49)}, nguncalake \VUgras{jala} menyang
tengahe \VUgras{kali} \textsuperscript{(50)}, oleh iwak kang luwih akeh.

%%K t08-c2-10-1 p
10. --- Mangsa gugur iku mangsane mlaku-mlaku kang becik : mlaku sikil,
\VUgras{nunggang jaran} \textsuperscript{(52)}, \VUgras{numpak pit}
\textsuperscript{(53)}, \VUgras{numpak montor} \textsuperscript{(54)}.
\VUgras{Dalan} \textsuperscript{(51)} padha resik lan wong nganggo
dina-dina endah kang pungkasan, amarga iki \VUgras{manuk sriti}
\textsuperscript{(56)} wis padha nglumpuk ing payon arep mabur sakayange
menyang negara kang luwih panas. Godhonge \VUgras{wit kenari}
\textsuperscript{(57)} tuwa wiwit kuning, \VUgras{woh kenari} padha tiba,
lan \VUgras{wit} dhuwur kang jejer-jejer kaya \VUgras{karangan kembang}
\textsuperscript{(58)} ing \VUgras{punthuk} \textsuperscript{(59)}
sedhela maneh bakal ora ana godhonge.

%%K t08-c2-11-1 p
11. --- \VUgras{Srengenge} \textsuperscript{(60)} munggah luwih kasep,
dinane saya cendhak, hawa adhem kang rada nyocog wiwit krasa ing
\VUgras{esuk,} lan iki pepanthane \VUgras{manuk bekasin}
\textsuperscript{(61)} wis padha ninggal iklim kita kang ora nampa
dhayoh.

\VUsaut{6.0mm}
\VUfilet{22mm}
